\subsection{向量丛的解析截面的限制}

设 X 是局部有限维的复解析流形，E 是以 X 为基的解析向量丛（VAR, p. 35 与 VAR, p. 70）。用 \(\mathcal{S}^{\omega}(\mathrm{X};\mathrm{E})\) 表示 E 在 X 上的解析截面所成的向量空间，赋以紧收敛拓扑。

引理 5. --- 设 \(X_{0}\) 是 X 的底实解析流形，而 \(E_{0}\) 是以 \(X_{0}\) 为基的、作为 E 的底的复向量丛。

\begin{enumerate}
\def\labelenumi{\alph{enumi})}
\item
  向量子空间 \(\mathcal{S}^{\omega}(\mathrm{X};\mathrm{E})\)（它包含于 \(\mathcal{S}^{0}(\mathrm{X}_{0};\mathrm{E}_{0})\) 中）是闭的，且从 \(\mathcal{S}^{\omega}(\mathrm{X};\mathrm{E})\) 到 \(\mathcal{S}^{0}(\mathrm{X}_{0};\mathrm{E}_{0})\) 中的单射是连续的且是严格的；
\item
  从 \(\mathcal{S}^{\omega}(\mathrm{X};\mathrm{E})\) 到 \(\mathcal{S}^{1}(\mathrm{X}_{0};\mathrm{E}_{0})\) 中的典范单射是连续的。
\end{enumerate}

首先假设存在整数 \(n \geqslant 0\) 与巴拿赫空间 F，使得 X 是 \(C^{n}\) 的开子集，且 E 是平凡向量丛 \(X \times F\)。在这种情形，拓扑向量空间 \(\mathcal{S}^{\omega}(X;E)\)、\(\mathcal{S}^{0}(X_{0};E_{0})\) 与 \(\mathcal{S}^{1}(X_{0};E_{0})\) 分别同构于 \(\mathcal{C}^{\omega}(X;F)\)、\(\mathcal{C}^{0}(X_{0};F)\) 与 \(\mathcal{C}^{1}(X_{0};F)\)，且空间 \(\mathcal{S}^{0}(X_{0};E_{0})\) 是可度量的（TG, X, p. 20, prop. 1 的 cor.）。引理于是由 VAR, 3.3.2, p. 28 得出。

设 \(\mathfrak{F}\) 是 \(\mathcal{S}^{\omega}(\mathrm{X};\mathrm{E})\) 上的一个滤子，它在空间 \(\mathcal{S}^{0}(\mathrm{X}_{0};\mathrm{E}_{0})\) 中收敛于极限 \(f\)。我们必须证明 \(f\) 属于 \(\mathcal{S}^{\omega}(\mathrm{X};\mathrm{E})\)，且滤子 \(\mathfrak{F}\) 收敛于 \(f\)（在空间 \(\mathcal{S}^{1}(\mathrm{X}_{0};\mathrm{E}_{0})\) 中）。这个命题是局部性质的，因而由证明的第一部分得出。

命题 11. --- 假设向量丛 E 是局部有限秩的。设 Y 是 X 的相对紧开子集。限制映射 \(f \mapsto f|Y\)（它从 \(\mathcal{S}^{\omega}(X;E)\) 到 \(\mathcal{S}^{\omega}(Y;E)\) 中）是紧的。

用引理 5 的记号，我们有交换图

\[
\begin{array}{c c c} \mathcal {S} ^ {\omega} (\mathrm{X}; \mathrm{E}) & \stackrel {{i}} {{\longrightarrow}} & \mathcal {S} ^ {1} (\mathrm{X} _ {0}; \mathrm{E} _ {0}) \\ \Big \downarrow^ {u} & & \Big \downarrow^ {v} \\ \mathcal {S} ^ {\omega} (\mathrm{Y}; \mathrm{E}) & \stackrel {{j}} {{\longrightarrow}} & \mathcal {S} ^ {0} (\mathrm{Y} _ {0}; \mathrm{E} _ {0}) \end{array}\tag{3}
\]

其中 i 与 j 是典范单射，而 u 与 v 是限制映射。映射 i 是连续的（引理 5, b)），映射 v 是紧的（Prop. 10）。由于典范单射 j 是连续的、严格的且有闭像（引理 5, a)），映射 u 是紧的（III, p. 3 的注记 5）。

推论. --- 设 X 是紧复解析流形，E 是 X 上的解析向量丛，且是局部有限秩的。向量空间 \(\mathcal{S}^{\omega}(X;E)\) 是有限维的。

解析流形 X 是紧的，因而是局部有限维的。由 prop. 11，\(\mathcal{S}^{\omega}(\mathrm{X};\mathrm{E})\) 的恒等映射是紧线性映射。这蕴含空间 \(\mathcal{S}^{\omega}(\mathrm{X};\mathrm{E})\) 是有限维的（III, p. 2 的注记 3）。

